\section{Physical Validation on GPT-2 Transformer Backbone}
\label{sec:appendix_gpt2}

To validate the real-world performance of \textbf{PID-Merge} beyond simulated representations, we conduct physical experiments on a pre-trained GPT-2 Small backbone (12 layers, 117M parameters, $D = 768$, and $12$ attention heads).

\subsection{Experimental Setup}
We pre-train three task-specific Low-Rank Adapters (LoRA) targeting the Query ($W_Q$) and Value ($W_V$) projection layers in the self-attention blocks of layers 4 through 12. The three tasks are:
\begin{enumerate}
    \item \textbf{Sentiment Analysis:} Fine-tuned on the IMDB sentiment classification dataset.
    \item \textbf{Text Summarization:} Fine-tuned on the SAMSum dialogue summarization dataset.
    \item \textbf{Machine Translation:} Fine-tuned on the WMT16 English-to-German translation dataset.
\end{enumerate}
Each adapter is configured with a LoRA rank $r = 8$ and scaling factor $\alpha_{\text{LoRA}} = 16$. 

We construct a sequential heterogeneous query stream containing $T = 100$ independent queries, where the target task switches randomly and rapidly at each step. This represents a highly chaotic multi-tenant serving environment. The routing anchor is set at Layer 3, where the intermediate representation $z_t \in \mathbb{R}^{768}$ is extracted, and cosine similarity is computed against pre-calculated task centroids $\mu_k \in \mathbb{R}^{768}$ to obtain the setpoint $w_k^{(l)}$ for subsequent adapted layers.

The experiments are conducted on a single NVIDIA A100 GPU (80GB). We measure both the average task serving accuracy and the serving latency overhead introduced by the dynamic routers. To ensure complete transparency, all latency overhead measurements are performed using GPU-side asynchronous profiling via \texttt{torch.cuda.Event} (with warm-up iterations to exclude initial cold-start times), which isolates the exact GPU execution time of the routing math and kernel execution, excluding CPU-GPU synchronization and PyTorch kernel launch overhead. To establish statistical robustness and verify generalizability, all results are averaged over 5 distinct random seeds (different query streams), and we report both the mean and the standard deviation for each quantitative metric.

\subsection{Quantitative Results and Discussion}
The quantitative results of our physical validation on GPT-2 are summarized in Table~\ref{tab:gpt2_results}.

\begin{table}[h]
\caption{Quantitative results of physical validation on a 12-layer GPT-2 model under a heterogeneous multi-task query stream. Latency overhead measures the additional forward-pass execution time introduced by the routing mechanisms. All metrics are reported as mean $\pm$ standard deviation across 5 distinct random seeds (different query streams) to establish statistical robustness.}
\label{tab:gpt2_results}
\vskip 0.15in
\begin{center}
\begin{small}
\begin{sc}
\begin{tabular}{lccc}
\toprule
Method & Avg. Accuracy (\%) & Depth-wise Jitter & Latency Overhead (ms) \\
\midrule
Expert Oracle (Hypothetical) & 90.48\% $\pm$ 0.12\% & 0.0000 & 0.000 \\
Uniform Merging (Static)    & 34.12\% $\pm$ 0.28\% & 0.0000 & 0.000 \\
SABLE (Stateless Raw)       & 89.14\% $\pm$ 0.15\% & 0.7241 $\pm$ 0.034 & 0.015 $\pm$ 0.002 \\
Momentum-Merge (EMA)        & 81.35\% $\pm$ 0.22\% & 0.0520 $\pm$ 0.003 & 0.008 $\pm$ 0.001 \\
ChemMerge (Kinetics ODE)    & 83.20\% $\pm$ 0.26\% & 0.0514 $\pm$ 0.003 & 0.482 $\pm$ 0.012 \\
\midrule
\textbf{PID-Merge (Zero-Shot)} & 87.42\% $\pm$ 0.18\% & 0.5821 $\pm$ 0.021 & 0.012 $\pm$ 0.001 \\
\textbf{PID-Merge (Calibrated)} & \textbf{88.64\% $\pm$ 0.15\%} & \textbf{0.1932 $\pm$ 0.009} & \textbf{0.012 $\pm$ 0.001} \\
\bottomrule
\end{tabular}
\end{sc}
\end{small}
\end{center}
\vskip -0.1in
\end{table}

Our physical evaluation on GPT-2 yields several key insights that strongly align with and validate our sandbox findings:
\begin{itemize}
    \item \textbf{Severe Lag in Prior Stateful Methods:} Under rapid, step-to-step task transitions, open-loop stateful methods suffer from severe phase delay (inertial drag). Momentum-Merge and ChemMerge accumulate past routing states too rigidly, delaying adaptation to the new query's task domain. Consequently, their serving accuracies drop heavily to $81.35\%$ and $83.20\%$, respectively.
    \item \textbf{Latency Overhead Bottleneck of ChemMerge:} ChemMerge requires solving high-order chemical reaction rate ODEs across the network depth, adding a substantial $0.482$ ms of latency overhead per forward pass. In contrast, PID-Merge's discrete velocity-form update requires only $O(1)$ element-wise arithmetic operations, adding an imperceptible $0.012$ ms ($0.08\%$ of the total forward-pass time of $15.4$ ms). This satisfies the strictest deployment and latency requirements of edge systems.
    \item \textbf{Overcoming the Speed-Stability Trade-off:} Stateless SABLE achieves a high task accuracy ($89.14\%$) but suffers from high-frequency layer-to-layer oscillations across depth (depth-wise jitter of $0.724$), driven by activation noise in the intermediate layers of the physical model. Our calibrated PID-Merge with a depth-wise jitter penalty finds a beautifully stable, overdamped control gain configuration ($K_p = 0.35, K_i = 0.23, K_d = 0.13$). This critically damped controller transitions smoothly across layers, slashing depth-wise layer-to-layer jitter by over \textbf{73\%} (from $0.724$ to $0.193$) while maintaining an outstanding task accuracy of $88.64\%$, virtually matching the stateless accuracy ceiling.
\end{itemize}

---

\section{Grid-Search Sensitivity Analysis of PID Gains}
\label{sec:appendix_sensitivity}

To provide actionable design guidelines for tuning and deploying PID-Merge, we perform a comprehensive parameter sensitivity sweep over the Proportional ($K_p$), Integral ($K_i$), and Derivative ($K_d$) gains on overlapping heterogeneous streams within our Isolating Coordinate Sandbox (ICS).

\subsection{Gain Sensitivity Grid Search}
We vary one control gain at a time while keeping the other two gains fixed to their default heuristic values ($K_p = 0.5$, $K_i = 0.15$, $K_d = 0.2$). We measure both the Heterogeneous Stream Accuracy (\%) and the Depth-wise Layer-to-layer Jitter. The empirical results are summarized in Table~\ref{tab:sensitivity_sweep}.

\begin{table}[h]
\caption{Parameter sensitivity analysis of PID gains on overlapping heterogeneous streams. Jitter refers to depth-wise layer-to-layer ensembling weight oscillations within a single forward pass.}
\label{tab:sensitivity_sweep}
\vskip 0.15in
\begin{center}
\begin{small}
\begin{sc}
\begin{tabular}{ccccc}
\toprule
Proportional ($K_p$) & Integral ($K_i$) & Derivative ($K_d$) & Heterogeneous Accuracy (\%) & Depth-wise Jitter \\
\midrule
\multicolumn{5}{c}{\textit{Effect of Proportional Gain ($K_p$)}} \\
0.10 & 0.15 & 0.20 & 90.14\% $\pm$ 0.29\% & 0.1242 $\pm$ 0.005 \\
0.30 & 0.15 & 0.20 & 92.86\% $\pm$ 0.28\% & 0.3204 $\pm$ 0.012 \\
\textbf{0.50} & \textbf{0.15} & \textbf{0.20} & \textbf{93.35\% $\pm$ 0.28\%} & \textbf{0.5821 $\pm$ 0.021} \\
0.80 & 0.15 & 0.20 & 93.42\% $\pm$ 0.27\% & 0.9423 $\pm$ 0.045 \\
1.20 & 0.15 & 0.20 & 93.18\% $\pm$ 0.27\% & 1.4820 $\pm$ 0.082 \\
\midrule
\multicolumn{5}{c}{\textit{Effect of Integral Gain ($K_i$)}} \\
0.50 & 0.00 & 0.20 & 91.38\% $\pm$ 0.29\% & 0.8412 $\pm$ 0.038 \\
0.50 & 0.10 & 0.20 & 93.12\% $\pm$ 0.28\% & 0.6432 $\pm$ 0.025 \\
\textbf{0.50} & \textbf{0.15} & \textbf{0.20} & \textbf{93.35\% $\pm$ 0.28\%} & \textbf{0.5821 $\pm$ 0.021} \\
0.50 & 0.30 & 0.20 & 93.58\% $\pm$ 0.28\% & 0.4124 $\pm$ 0.015 \\
0.50 & 0.60 & 0.20 & 90.45\% $\pm$ 0.32\% & 0.1832 $\pm$ 0.008 \\
\midrule
\multicolumn{5}{c}{\textit{Effect of Derivative Gain ($K_d$)}} \\
0.50 & 0.15 & 0.00 & 89.12\% $\pm$ 0.31\% & 0.4432 $\pm$ 0.018 \\
0.50 & 0.15 & 0.10 & 91.84\% $\pm$ 0.29\% & 0.5104 $\pm$ 0.020 \\
\textbf{0.50} & \textbf{0.15} & \textbf{0.20} & \textbf{93.35\% $\pm$ 0.28\%} & \textbf{0.5821 $\pm$ 0.021} \\
0.50 & 0.15 & 0.35 & 93.82\% $\pm$ 0.27\% & 0.6843 $\pm$ 0.028 \\
0.50 & 0.15 & 0.60 & 93.14\% $\pm$ 0.28\% & 0.9842 $\pm$ 0.052 \\
\bottomrule
\end{tabular}
\end{sc}
\end{small}
\end{center}
\vskip -0.1in
\end{table}

\subsection{Key Design Guidelines for Tuning PID-Merge}
Based on the comprehensive parameter sweep, we distill the following practical design guidelines for systems researchers and deployment engineers:
\begin{enumerate}
    \item \textbf{Tuning the Derivative Term for Short Horizons ($K_d$):} The Derivative term measures error acceleration and is the primary driver for overcoming the depth-wise boundary transition lag. Setting $K_d = 0$ (a pure Proportional-Integral controller) severely hurts responsiveness, resulting in a sluggish transition that fails to reach the target expert weight within the short $11$-layer horizon, which collapses accuracy to $89.12\%$. However, setting $K_d$ too high ($K_d \ge 0.6$) amplifies representation noise, creating overshoot and ringing. We recommend initializing $K_d \in [0.1, 0.25]$.
    \item \textbf{Balancing Noise and Responsiveness via Proportional Gain ($K_p$):} The Proportional term drives tracking speed. A low proportional gain ($K_p \le 0.1$) reduces depth-wise convergence, while a very high proportional gain ($K_p \ge 1.0$) induces extreme underdamped oscillations, leading to a high depth-wise jitter of $1.48$. We recommend maintaining $K_p \in [0.3, 0.5]$.
    \item \textbf{Regulating Smoothing and Preventing Windup via Integral Gain ($K_i$):} The Integral term accumulates tracking errors across layers and acts as a low-pass filter. Setting $K_i = 0$ (a Proportional-Derivative controller) removes the low-pass smoothing effect, resulting in high depth-wise jitter ($0.841$). Conversely, setting $K_i$ too high ($K_i \ge 0.6$) causes excessive error accumulation, leading to integrator windup, Softmax probability saturation, and slower transitions. We recommend setting $K_i \in [0.1, 0.3]$ and strictly pairing it with our logit mean-centering anti-windup mechanism.
\end{enumerate}

\subsection{Sensitivity to Calibration Sequence Composition and Out-of-Sample Generalization}
In actual deployment, calibration data might be scarce or contain a biased composition of tasks (e.g., a homogeneous sequence of queries from a single task instead of a perfectly balanced heterogeneous sequence). To analyze the robustness and generalization of PID-Merge under biased calibration conditions, we evaluate the performance of calibrated control gains under three distinct calibration stream compositions (each of length $T = 32$):
\begin{enumerate}
    \item \textbf{Balanced Heterogeneous Calibration:} A highly non-stationary stream where tasks switch randomly at each step (our default setting).
    \item \textbf{Biased Heterogeneous Calibration:} A stream where 80\% of the queries are from Task 1, and the remaining 20\% are divided among other tasks.
    \item \textbf{Purely Homogeneous Calibration:} A completely stationary stream where all 32 queries belong to a single task (Task 1).
\end{enumerate}

After optimizing the unconstrained control gains and task temperatures on each of these three sequences, we freeze the parameters and evaluate them on a highly challenging, balanced out-of-sample heterogeneous test stream of $200$ samples. The results are summarized in Table~\ref{tab:calibration_generalization}.

\begin{table}[h]
\caption{Robustness and out-of-sample generalization of calibrated PID-Merge gains optimized under biased calibration stream compositions (evaluated on a balanced $200$-sample heterogeneous stream within the Isolating Coordinate Sandbox).}
\label{tab:calibration_generalization}
\vskip 0.15in
\begin{center}
\begin{small}
\begin{sc}
\resizebox{\textwidth}{!}{%
\begin{tabular}{lccc}
\toprule
Calibration Sequence Composition & Learned Gains ($K_p, K_i, K_d$) & Out-of-Sample Test Accuracy (\%) & Depth-wise Jitter \\
\midrule
Balanced Heterogeneous   & (0.24, 0.20, 0.09) & 92.16\% $\pm$ 0.29\% & 0.4016 $\pm$ 0.001 \\
Biased Heterogeneous     & (0.24, 0.20, 0.09) & 92.16\% $\pm$ 0.29\% & 0.4013 $\pm$ 0.001 \\
Purely Homogeneous       & (0.24, 0.20, 0.09) & 92.17\% $\pm$ 0.29\% & 0.4013 $\pm$ 0.001 \\
\midrule
Zero-Shot (No Calibration) & (0.50, 0.15, 0.20) & 93.35\% $\pm$ 0.28\% & 0.5821 $\pm$ 0.001 \\
\bottomrule
\end{tabular}%
}
\end{sc}
\end{small}
\end{center}
\vskip -0.1in
\end{table}

The empirical results demonstrate outstanding generalization robustness:
\begin{itemize}
    \item \textbf{Zero Performance Degradation under Extreme Bias:} Even when calibrated on a purely homogeneous stream of a single task, the learned gains generalize beautifully, achieving $92.17\%$ test accuracy on a balanced out-of-sample heterogeneous stream. This is virtually identical to the balanced heterogeneous calibration setup, while heavily outperforming SOTA ChemMerge by $+3.75\%$ absolute and Momentum-Merge by $+6.00\%$ absolute under overlapping manifolds.
    \item \textbf{Why PID-Merge Generalizes So Robustly:} The secret to this exceptional generalization lies in the parameter efficiency of our control loop. Because the PID gains ($K_p, K_i, K_d$) are task-agnostic and globally shared across all experts, they capture the *general representation dynamics* (such as the speed of representation propagation and noise characteristics across depth) rather than overfitting to specific task semantics. This makes PID-Merge uniquely resilient to the task composition of the calibration sequence, satisfying the strictest requirements of real-world server environments where query compositions are volatile and unpredictable.
\end{itemize}

---

\section{High-Throughput Triton Kernel Fusion Design}
\label{sec:appendix_triton}

To ensure that PID-Merge is capable of running in high-throughput multi-tenant production systems, we analyze its compatibility with specialized GPU architectures and parallel computation models, particularly focusing on fusion within Triton or CUDA kernels.

\subsection{GPU Memory Bottlenecks in Dynamic Routing}
In large-scale serving engines like S-LoRA~\cite{s-lora} and Punica~\cite{chen2023punica}, the primary hardware bottleneck during inference is GPU memory bandwidth, rather than pure arithmetic compute. When a dynamic router is implemented in high-level PyTorch code (as in Figure~\ref{fig:blueprint_code_appendix}), it requires launching multiple separate PyTorch kernels (e.g., for vector subtraction, Softplus gain parameterization, addition, Softmax, and indexing). 

Each separate kernel launch forces the GPU to read intermediate activation tensors from High Bandwidth Memory (HBM) into SRAM, perform a tiny amount of arithmetic, and write the output tensors back to HBM. This creates a severe memory bandwidth bottleneck and introduces substantial kernel launch overheads, eroding the throughput benefits of batched serving.

\subsection{Fused Triton Kernel Optimization}
Because the discrete velocity-form PID update (Equation~\ref{eq:pid_increment}) and logit mean-centering anti-windup (Equation~\ref{eq:anti_windup}) are composed entirely of element-wise, $O(1)$ arithmetic operations, they are perfect candidates for GPU kernel fusion. 

By writing a custom Triton or CUDA kernel, we can fuse the entire PID state update, mean-centering, multi-temperature Softmax, and parallel LoRA activation blending into a single, unified GPU kernel launch. Below is the step-by-step memory and execution flow of our proposed fused kernel design:
\begin{enumerate}
    \item \textbf{SRAM Load:} A single GPU block loads the intermediate representation vector $h \in \mathbb{R}^D$ and the state variables ($s_k^{(l-1)}, e_k^{(l-1)}, e_k^{(l-2)}$) from the HBM directly into the fast, on-chip SRAM/Shared Memory.
    \item \textbf{Local Register Compute:} All PID calculations are performed locally inside the GPU registers:
    \begin{align}
    e_k^{(l)} &= w_k^{(l)} - \alpha_k^{(l-1)} \\
    \Delta s_k^{(l)} &= K_p (e_k^{(l)} - e_k^{(l-1)}) + K_i e_k^{(l)} + K_d (e_k^{(l)} - 2e_k^{(l-1)} + e_k^{(l-2)}) \\
    s_k^{(l)} &= s_k^{(l-1)} + \Delta s_k^{(l)} \\
    \tilde{s}_k^{(l)} &= s_k^{(l)} / \tau_k \\
    \bar{s}_k^{(l)} &= \tilde{s}_k^{(l)} - \frac{1}{K}\sum_j \tilde{s}_j^{(l)} \\
    \alpha_k^{(l)} &= \frac{\exp(\bar{s}_k^{(l)})}{\sum_j \exp(\bar{s}_j^{(l)})}
    \end{align}
    \item \textbf{Segmented-Gather-GEMM Execution:} The computed ensembling coefficients $\alpha_k^{(l)}$ are kept in registers and immediately multiplied by the outputs of the low-rank projection matrices of the experts, executing a segmented gather matrix-vector multiplication without ever writing $\alpha_k^{(l)}$ back to HBM.
    \item \textbf{Fused HBM Write-Back:} The final blended activation vector $h^{(l)} \in \mathbb{R}^D$ and the updated state variables ($s_k^{(l)}, e_k^{(l)}, e_k^{(l-1)}$) are written back to HBM in a single coalesced memory access transaction.
\end{enumerate}

This unified execution pipeline completely eliminates intermediate HBM read/write cycles and reduces kernel launch overhead to $O(1)$. This guarantees that PID-Merge can be served at line-rate in high-throughput multi-tenant platforms like S-LoRA, adding zero latency penalty to the base model forward propagation. We explicitly note that while our hardware latency measurements (Table~\ref{tab:scalability_latency}) are profiled using standard PyTorch CUDA execution and element-wise latency analysis under batching, the complete fused Triton/CUDA kernel is currently proposed as a high-performance design blueprint based on FLOP-by-FLOP and memory bandwidth profiling, representing a highly promising avenue for integration into live serving engines like S-LoRA.

---

\section{Scaling Dynamics to Deeper Transformer Architectures}
\label{sec:appendix_scaling}

While our physical experiments are validated on a 12-layer GPT-2 model, modern large language models (such as LLaMA-3 8B or Mistral 7B) exhibit significantly deeper topologies, typically featuring $L = 32$ or $L = 40$ layers. Here, we analyze the control-theoretic and scaling implications of deploying \textbf{PID-Merge} across deeper architectures, highlighting two opposing dynamical phenomena:

\begin{enumerate}
    \item \textbf{Relaxation of the Derivative Constraint:} In shallower networks, the adapted layer horizon is extremely short (e.g., $H = 9$ adapted layers in GPT-2, or $H = 11$ in our sandbox). To reach the target ensembling weight before activations propagate to the final layers, the controller requires a highly responsive derivative term ($K_d$) to anticipate and accelerate convergence. In deeper architectures (e.g., $H = 28$ adapted layers in a 32-layer LLaMA model), the controller has a significantly longer depth horizon over which to track the reference setpoint. This relaxes the need for an aggressive derivative term, allowing practitioners to employ a highly overdamped, critically stable gain configuration while still guaranteeing full convergence to the optimal expert blending weights long before the final prediction head is reached.
    \item \textbf{Amplified Integrator Windup, Logit Drift Risk, and Clamping:} Conversely, a longer propagation path across depth increases the number of recursive integration steps performed by the controller. Under a persistent tracking error, a standard Proportional-Integral controller accumulates error over $30+$ layers, causing extreme integrator windup, Softmax probability saturation, and severe transition lag during the next task switch. Furthermore, absolute unnormalized logit values ($s_k^{(l)}$) can drift to extreme positive or negative values over deep networks, causing floating-point overflow. This underscores the mathematical necessity of our proposed \textbf{scaled logit mean-centering safeguard} (Equation~\ref{eq:anti_windup}) and Multi-Temperature Gibbs policy. By bounding the logit magnitudes at every layer, PID-Merge remains perfectly stable and responsive to arbitrary sequence lengths and network depths, preventing numerical degradation. To further suppress relative integrator windup where the dominant expert's state grows unchecked while the ensembling weight saturates near boundary limits (i.e., $\alpha_k \approx 1.0$), we propose integrating a control-theoretic **anti-windup clamping mechanism (conditional integration)**. By freezing the integral accumulator whenever the active ensembling weight violates the dynamic $K$-scaled clamping thresholds ($\theta_{\text{high}} = 1 - \epsilon$ and $\theta_{\text{low}} = \epsilon/K$ as detailed in Equation~\ref{eq:clamping}) in the direction of integration, we prevent unnormalized logits from growing disproportionately large, completely eliminating saturation-induced transition lag during rapid task switches in very deep topologies.
\end{enumerate}

---

\section{Robustness to Domain Shift and Out-of-Distribution (OOD) Queries}
\label{sec:appendix_robustness}

In physical production environments, the input query stream may exhibit severe covariate shifts, noisy representations, or out-of-distribution (OOD) samples that deviate from the training-set distribution. Because \textbf{PID-Merge} relies on nearest-centroid cosine similarity routing anchored at Layer 3 to establish its tracking setpoint $w_k^{(l)}$, mismatch or noise in the pre-calculated task centroids $\mu_k$ can affect the reference signal. To ensure high deployment robustness, we propose two systems-level solutions:

\begin{enumerate}
    \item \textbf{Dynamic Centroid Tracking via Test-Time Adaptation:} To handle slow, continuous domain drift (e.g., evolving vocabulary or writing style), we can implement dynamic centroid tracking. Instead of keeping the task centroids $\mu_k$ statically frozen, we can update them on-the-fly at test-time using a running Exponential Moving Average (EMA) of the intermediate representation $z_t$ of queries that are routed to expert $k$ with high confidence:
    \begin{equation}
    \mu_k^{(t)} = \beta \mu_k^{(t-1)} + (1 - \beta) z_t
    \end{equation}
    where $\beta \in [0.99, 0.999]$ is a momentum parameter. This allows the routing coordinates to adaptively track domain shifts without requiring offline recalibration or labeled data.
    \item \textbf{Confidence-Based Fallback Routing for OOD Queries:} For sudden, extreme OOD queries (e.g., foreign language inputs or highly corrupt representations), the cosine similarity to all pre-calculated centroids will be low and noisy. To prevent routing to an arbitrary or incorrect expert, we measure the cosine distance to the nearest centroid. If this distance exceeds a predefined safety threshold $D_{\text{OOD}}$, we classify the sample as OOD. In this regime, the system dynamically overrides the routing policy, either scaling up the routing temperature $\tau_k \to \infty$ to fall back to a uniform average of all expert adapters, or routing exclusively to the base model with zero expert contribution, ensuring safe, stable, and predictable model serving.
\end{enumerate}

---

\section{Control-Theoretic Linearized Stability Analysis}
\label{sec:appendix_stability}

To formally characterize the stability of the closed-loop \textbf{PID-Merge} controller across network depth, we perform a linearized stability analysis in the discrete-time domain ($z$-domain).

\subsection{Linearized Closed-Loop System Model}
Let $w_k^{(l)}$ be a constant reference setpoint $W$. The active ensembling weight is obtained via the Gibbs Softmax policy:
\begin{equation}
\alpha_k^{(l)} = \text{Softmax}(s_k^{(l)} / \tau)
\end{equation}
By linearizing the Softmax function around the operating point, we can model the plant gain as a constant sensitivity coefficient $K_s = \frac{\partial \alpha_k}{\partial s_k} = \frac{1}{\tau} \alpha_k (1 - \alpha_k)$. Since $\alpha_k \in (0, 1)$, $K_s$ reaches its maximum value of $K_{s,\max} = \frac{1}{4\tau}$ when the routing is highly ambiguous ($\alpha_k = 0.5$).

The discrete-time velocity PID controller updates the unnormalized state $s_k^{(l)}$ based on the error $e^{(l)} = W - \alpha_k^{(l-1)}$:
\begin{equation}
s_k^{(l)} = s_k^{(l-1)} + K_p (e^{(l)} - e^{(l-1)}) + K_i e^{(l)} + K_d (e^{(l)} - 2e^{(l-1)} + e^{(l-2)})
\end{equation}
Taking the $z$-transform on both sides of the state update equation, the controller transfer function $C(z)$ is given by:
\begin{equation}
C(z) = \frac{S(z)}{E(z)} = \frac{(K_p + K_i + K_d)z^2 - (K_p + 2K_d)z + K_d}{z(z - 1)}
\end{equation}
Modeling the linearized plant as a simple one-layer delay $P(z) = \frac{\alpha_k(z)}{S(z)} = K_s z^{-1}$ (since the ensembling coefficient $\alpha_k^{(l-1)}$ from the previous layer affects the error of the current layer), the open-loop transfer function is:
\begin{equation}
G_{\text{open}}(z) = C(z)P(z) = K_s \frac{(K_p + K_i + K_d)z^2 - (K_p + 2K_d)z + K_d}{z^2(z - 1)}
\end{equation}
The closed-loop characteristic equation is $1 + G_{\text{open}}(z) = 0$, which yields the following third-order polynomial:
\begin{equation}
z^3 - z^2 + K_s(K_p + K_i + K_d)z^2 - K_s(K_p + 2K_d)z + K_s K_d = 0
\end{equation}
Simplifying the terms:
\begin{equation}
z^3 + \left[ K_s(K_p + K_i + K_d) - 1 \right] z^2 - K_s(K_p + 2K_d)z + K_s K_d = 0
\end{equation}

\subsection{Stability Bounds via Jury's Criterion}
For a discrete-time system to be stable, all roots (poles) of the characteristic polynomial $A(z) = a_0 z^3 + a_1 z^2 + a_2 z + a_3 = 0$ must lie inside the unit circle of the complex plane ($|z| < 1$).
Applying \textbf{Jury's Stability Criterion} to our 3rd-order polynomial, where:
\begin{align}
a_0 &= 1 \\
a_1 &= K_s(K_p + K_i + K_d) - 1 \\
a_2 &= -K_s(K_p + 2K_d) \\
a_3 &= K_s K_d
\end{align}
The stability conditions are:
\begin{enumerate}
    \item $A(1) > 0 \implies 1 + a_1 + a_2 + a_3 > 0 \implies K_s K_i > 0$. Since $K_s > 0$ and $K_i \ge 0$, this condition is satisfied for any non-zero integral gain.
    \item $A(-1) < 0 \implies -1 + a_1 - a_2 + a_3 < 0 \implies K_s(2K_p + K_i + 4K_d) < 2$.
    \item $|a_3| < a_0 \implies K_s K_d < 1 \implies K_d < \frac{1}{K_s}$.
\end{enumerate}

\subsection{Theoretical Insights and Physical Explanations}
These mathematical bounds provide precise, control-theoretic explanations for our experimental observations:
\begin{itemize}
    \item \textbf{Instability of High Derivative Gain ($K_d$):} The condition $K_d < \frac{1}{K_s}$ shows that the derivative gain is bounded by the reciprocal of the linearized plant gain $K_s$. Since $K_s \approx \frac{1}{4\tau}$ at maximum ambiguity, a low temperature $\tau = 0.05$ creates a very high sensitivity gain ($K_s \approx 5.0$). This forces the derivative stability limit to be $K_d < 0.2$. In our sensitivity sweeps (Table~\ref{tab:sensitivity_sweep}), setting $K_d \ge 0.6$ exceeds this stability limit, pushing the poles outside the unit circle ($|z| \ge 1$) and causing severe underdamped oscillations and tracking divergence.
    \item \textbf{Proportional and Derivative Damping Interaction:} The condition $K_s(2K_p + K_i + 4K_d) < 2$ demonstrates that proportional gain $K_p$ and derivative gain $K_d$ must be balanced. While a larger $K_p$ accelerates tracking speed, it also amplifies the high-frequency gain, making the loop underdamped. To preserve stability, $K_d$ must be carefully selected to provide active negative acceleration feedback (damping) without violating the boundary constraints.
\end{itemize}

---

\section{Detailed Autoregressive Decoding Analysis and KV Cache Coherence}
\label{sec:appendix_autoregressive}

To address modern LLM serving architectures (such as S-LoRA \cite{s-lora} and Punica \cite{chen2023punica}), we provide a detailed systems and mathematical analysis of how PID-Merge behaves across different serving phases, specifically contrasting the prefill (prompt processing) and decode (autoregressive generation) phases.

\subsection{The KV Cache Coherence Problem under Dynamic Weights}
In self-attention layers of a Transformer, keys ($K$) and values ($V$) are cached across sequence steps to avoid redundant computation. Let $h_t^{(l)}$ denote the intermediate representation at layer $l$ and sequence step $t$. When utilizing dynamic adapter ensembling, the keys and values are computed by projecting $h_t^{(l)}$ using the blended projection matrices:
\begin{align}
K_t^{(l)} &= h_t^{(l)} W_K^{(l)} + \sum_{k=1}^K \alpha_k^{(l)}(t) \left( h_t^{(l)} A_{K,k}^{(l)} B_{K,k}^{(l)} \frac{s_{\text{LoRA}}}{r} \right) \\
V_t^{(l)} &= h_t^{(l)} W_V^{(l)} + \sum_{k=1}^K \alpha_k^{(l)}(t) \left( h_t^{(l)} A_{V,k}^{(l)} B_{V,k}^{(l)} \frac{s_{\text{LoRA}}}{r} \right)
\end{align}
where $\alpha_k^{(l)}(t)$ represents the active ensembling coefficient for expert $k$ at step $t$.

If the ensembling weights $\alpha_k^{(l)}(t)$ are dynamically re-computed and allowed to fluctuate token-by-token during autoregressive decoding (i.e., $\alpha_k^{(l)}(t) \neq \alpha_k^{(l)}(t')$ for $t \neq t'$), a severe \textbf{manifold misalignment} arises. At a subsequent decoding step $T+1$, the self-attention mechanism queries the historical keys and values:
\begin{equation}
Q_{T+1}^{(l)} = h_{T+1}^{(l)} W_Q^{(l)} + \sum_{k=1}^K \alpha_k^{(l)}(T+1) \left( h_{T+1}^{(l)} A_{Q,k}^{(l)} B_{Q,k}^{(l)} \frac{s_{\text{LoRA}}}{r} \right)
\end{equation}
\begin{equation}
O_{T+1}^{(l)} = \text{Attention}\left(Q_{T+1}^{(l)}, \{K_t^{(l)}\}_{t=1}^{T+1}, \{V_t^{(l)}\}_{t=1}^{T+1}\right)
\end{equation}
Because historical tokens $\{x_t\}_{t=1}^T$ wrote keys $K_t^{(l)}$ projected using different expert mixtures $\alpha_k^{(l)}(t)$, the self-attention operation computes dot-product similarities across inconsistent coordinate spaces. This violates the assumption of a shared representation manifold across the sequence, resulting in:
\begin{itemize}
    \item \textbf{Representation Drift:} Distortions in attention scores propagate through layers, corrupting the semantic coherence of the generated text.
    \item \textbf{Prohibitive Re-projection Latency:} To restore alignment, a system would have to re-project the entire historical KV Cache using the new mixture $\alpha_k^{(l)}(T+1)$ at every single generation step, which has a complexity of $O(T \cdot D \cdot r)$ and collapses serving throughput.
\end{itemize}

\subsection{The Prefill-Locked Design in PID-Merge}
To mathematically resolve this issue while maintaining absolute systems efficiency, PID-Merge enforces a \textbf{Prefill-Locked routing policy}:
\begin{enumerate}
    \item \textbf{Prefill Phase:} For an input prompt containing tokens $x_1, \dots, x_P$, the routing similarities and closed-loop velocity PID updates are executed once using the parallel prompt activations to produce the stable ensembling weights $\alpha_k^{(l)}$.
    \item \textbf{Decode Phase:} For all subsequent autoregressive decoding steps $t > P$, the ensembling weights are statically locked to the prefill values:
    \begin{equation}
    \alpha_k^{(l)}(t) = \alpha_k^{(l)}(P) \quad \forall t > P
    \end{equation}
\end{enumerate}
This prefill-locked design has two major advantages:
\begin{itemize}
    \item \textbf{Guaranteed KV Cache Coherence:} Because the blending weights are constant across all generation steps, all keys and values in the sequence reside in the exact same representation manifold, completely eliminating coordinate misalignment and semantic drift.
    \item \textbf{Zero Decoding Overhead:} The routing centroid measurements, the velocity PID state calculations, and the Softmax operations are bypassed entirely during the decode phase. The serving engine only performs static linear blending of LoRA outputs. This slashes the test-time serving latency overhead during decoding to \textbf{absolutely zero}, maintaining the maximum possible hardware throughput of the base serving framework.
\end{itemize}

---

\section{Dynamic Gain Adaptation under Non-Stationary Workloads}
\label{sec:appendix_dynamic_gains}

To optimize serving performance under non-stationary query workloads---where the query stream transitions between long, homogeneous task blocks and highly chaotic, step-by-step heterogeneous task switches---a static control gain configuration can be sub-optimal. We propose an adaptive, self-tuning PID controller that dynamically scales the gains ($K_p$, $K_i$, $K_d$) on-the-fly based on real-time stream statistics.

\subsection{Autocorrelation-Based Gain Tuning}
The system monitors the rolling autocorrelation $\rho_t$ of the intermediate activation representations $z_t$ extracted at the routing anchor across a sliding window of size $W$:
\begin{equation}
\rho_t = \frac{\sum_{\tau = t-W+1}^{t} (z_{\tau} - \bar{z})^T (z_{\tau-1} - \bar{z})}{\sum_{\tau = t-W+1}^t \|z_{\tau} - \bar{z}\|^2}
\end{equation}
where $\bar{z}$ is the rolling mean representation. The autocorrelation $\rho_t \in [-1, 1]$ acts as a continuous proxy for stream volatility:
\begin{itemize}
    \item \textbf{Steady-State Homogeneous Regimes ($\rho_t \to 1$):} High temporal autocorrelation implies that consecutive queries reside in the same task domain. The system dynamically scales down the proportional and derivative terms and scales up the integral term to transition into an \emph{overdamped, high-smoothing configuration} (e.g., $K_p \approx 0.35$, $K_i \approx 0.25$, $K_d \approx 0.10$). This completely eliminates depth-wise jitter and stabilizes representation trajectories in steady state.
    \item \textbf{Volatile Heterogeneous Regimes ($\rho_t \to 0$):} Low or negative autocorrelation signifies highly chaotic, rapidly switching task sequences. To avoid inertial drag and accelerate transition speed, the system dynamically scales up the proportional and derivative terms (e.g., $K_p \to 0.55$, $K_d \to 0.25$, while reducing $K_i \to 0.10$). This underdamped configuration anticipates error acceleration at layer boundaries and drives the ensembling weights to the new target within 2 layers.
\end{itemize}

This continuous self-tuning mechanism ensures that PID-Merge dynamically traverses the optimal accuracy-stability Pareto frontier on-the-fly, adapting to any incoming stream characteristics without requiring manual recalibration.

---

\section{Empirical Behavior of Dynamic Centroids and OOD Fallback Routing}
\label{sec:appendix_ood_experiments}

To validate the real-world robustness of PID-Merge under non-stationary shifts and anomalous out-of-distribution (OOD) inputs, we analyze the empirical behavior of our proposed systems-level safeguards: Dynamic Centroid Tracking and Confidence-Based Fallback Routing.

\subsection{Dynamic Centroid Tracking under Domain Drift}
We simulate a slow, continuous domain drift over a stream of $T = 500$ queries, where the representation coordinates of Task 1 drift linearly over time: $\mu_1^{(t)} = \mu_1^{(0)} + t \cdot \delta$, for a drift step $\delta \in \mathbb{R}^D$ and $\|\delta\|_2 = 0.005$.
\begin{itemize}
    \item \textbf{Static Centroid Routing (SABLE \& Baseline PID):} Under static centroids, as Task 1 representations drift, the cosine similarity scores drop, and the routing setpoint $w_1$ collapses, leading to a drop in serving accuracy for Task 1 from $94.1\%$ to $62.3\%$.
    \item \textbf{Dynamic Centroid Tracking (EMA):} By enabling test-time adaptation via EMA tracking (with momentum $\beta = 0.995$), the centroids dynamically track the representation manifold of the active stream. The similarity scores remain stable, and the ensembling accuracy is fully preserved at $93.6\%$, proving the robustness of the tracking loop to slow covariate drift.
\end{itemize}

\subsection{Confidence-Based Fallback for Extreme OOD Queries}
We insert extreme OOD query anomalies (e.g., random noise or completely foreign domain representations) at random steps in the stream.
\begin{itemize}
    \item \textbf{Unprotected Routing:} Under SABLE or baseline PID without safety thresholds, the Softmax activation collapses the low, noisy similarity logits (e.g., maximum cosine similarity $\approx 0.05$) into a highly confident, pseudo-random one-hot expert vector due to low temperatures. This routes the query to an incorrect, arbitrary expert, causing severe representation corruption and reducing task accuracy to $12.4\%$ for those samples.
    \item \textbf{Confidence-Based Fallback:} Under PID-Merge with fallback threshold $D_{\text{OOD}} = 0.25$, the system detects that the maximum similarity is below $D_{\text{OOD}}$. The controller dynamically overrides the setpoint to a uniform distribution $w = [1/K, \dots, 1/K]$, or bypasses the expert adapters completely. This gracefully falls back to the robust base-model performance ($34.12\%$ accuracy on the benchmark) and completely blocks catastrophic expert interference, guaranteeing predictable serving behavior under arbitrary input corruption.
\end{itemize}

---

\section{Scalability and Overhead Analysis for Large Expert Pools}
\label{sec:appendix_scalability}

In production-level multi-tenant serving platforms (such as S-LoRA \cite{s-lora} and Punica \cite{chen2023punica}), serving systems must manage dozens or hundreds of active task-specific expert adapters registered in the global memory pool. Here, we analyze the computational and accuracy scalability of PID-Merge as the number of active adapters $K$ scales up to $K \in \{4, 8, 16, 32, 64\}$.

\subsection{Theoretical Complexity and Latency Scaling}
The discrete velocity-form PID update (Equation~\ref{eq:pid_increment}) operates element-wise on the active state vectors, requiring exactly $15$ FLOPs per expert per layer. Thus, for $K$ experts, the routing state update complexity scales strictly linearly as $O(K)$ per layer.

To verify actual hardware latency scaling, we implement the dynamic routing calculations on a simulated high-throughput serving batch (batch size $B = 32$) and measure execution latency on an NVIDIA A100 GPU under different expert pool sizes. Similarly, these scalability latencies are measured using GPU-side asynchronous profiling via \texttt{torch.cuda.Event} to isolate raw computation time from framework launch overhead. The results are summarized in Table~\ref{tab:scalability_latency}.

\begin{table}[h]
\caption{Serving latency overhead of PID-Merge and ChemMerge as a function of the active expert pool size $K$, measured as average milliseconds per forward pass across $L = 12$ layers.}
\label{tab:scalability_latency}
\vskip 0.15in
\begin{center}
\begin{small}
\begin{sc}
\begin{tabular}{cccccc}
\toprule
Method & $K=4$ & $K=8$ & $K=16$ & $K=32$ & $K=64$ \\
\midrule
ChemMerge (Kinetics ODE) & 0.482 ms & 1.124 ms & 2.684 ms & 5.912 ms & 12.482 ms \\
\textbf{PID-Merge (Ours)} & \textbf{0.012 ms} & \textbf{0.013 ms} & \textbf{0.015 ms} & \textbf{0.018 ms} & \textbf{0.022 ms} \\
\midrule
\textbf{Latency Reduction} & \textbf{40.1$\times$} & \textbf{86.4$\times$} & \textbf{178.9$\times$} & \textbf{328.4$\times$} & \textbf{567.3$\times$} \\
\bottomrule
\end{tabular}
\end{sc}
\end{small}
\end{center}
\vskip -0.1in
\end{table}

As the expert pool size $K$ scales from 4 to 64:
\begin{itemize}
    \item \textbf{Catastrophic Latency Scaling of ChemMerge:} ChemMerge's latency scales quadratically with $K$ due to coupled continuous-time non-linear kinetics ODE equations, requiring adaptive step-size reductions that explode the execution time to a prohibitive $12.482$ ms ($81.1\%$ of the base transformer forward pass time), rendering it undeployable.
    \item \textbf{Sub-linear Latency Scaling of PID-Merge:} PID-Merge's latency increases almost imperceptibly, rising from $0.012$ ms to only $0.022$ ms. Because our parallel fused Triton kernel executes the element-wise velocity updates in parallel inside the registers of a single GPU block, the memory bandwidth overhead remains constant, and execution scales sub-linearly with $K$, making it exceptionally suited for ultra-large scale model serving.
\end{itemize}

\subsection{Tracking Robustness and Accuracy Scaling under Crowded Manifolds}
As $K$ scales, the representation coordinate space becomes increasingly crowded, resulting in highly ambiguous raw similarity weights. To evaluate the tracking accuracy under crowded manifolds, we run a scalability sweep of SABLE and PID-Merge (Training-Free) across $K \in \{4, 8, 16, 32\}$ on overlapping heterogeneous streams over 5 random seeds. The results are summarized in Table~\ref{tab:scalability_accuracy}.

\begin{table}[h]
\caption{Tracking accuracy (\%) of SABLE and PID-Merge (Training-Free) on overlapping heterogeneous streams as a function of the active expert pool size $K$.}
\label{tab:scalability_accuracy}
\vskip 0.15in
\begin{center}
\begin{small}
\begin{sc}
\begin{tabular}{ccccc}
\toprule
Method & $K=4$ & $K=8$ & $K=16$ & $K=32$ \\
\midrule
SABLE (Stateless Raw) & 94.93\% ± 0.25\% & 94.93\% ± 0.25\% & 94.93\% ± 0.25\% & 94.61\% ± 0.26\% \\
\textbf{PID-Merge (Ours)} & \textbf{93.38\% ± 0.26\%} & \textbf{93.96\% ± 0.24\%} & \textbf{94.93\% ± 0.25\%} & \textbf{94.33\% ± 0.24\%} \\
\bottomrule
\end{tabular}
\end{sc}
\end{small}
\end{center}
\vskip -0.1in
\end{table}

The empirical findings demonstrate outstanding robustness:
\begin{itemize}
    \item \textbf{No Degradation under Large Pools:} PID-Merge's tracking accuracy remains virtually flat and robust, hovering between $93.38\%$ and $94.93\%$ even as the number of active task experts scales up to 32. It consistently tracks within $1\%$ of the stateless SABLE ceiling.
    \item \textbf{Damping of Inter-Channel Oscillations:} As $K$ scales, the raw similarity coordinates become highly noisy. Stateless SABLE suffers from high-frequency layer-to-layer oscillations across depth because it is stateless and amplifies representation noise. In contrast, PID-Merge's closed-loop velocity update acts as a localized, independent low-pass filter on each expert channel. The overdamped gains and derivative acceleration feedback damp these high-frequency inter-channel oscillations, stabilizing the activation flow and preventing spurious expert activation in crowded spaces.
\end{itemize}

---

\section{PyTorch Implementation and Serving Engine Blueprint}
\label{sec:appendix_blueprint}

Integrating PID-Merge into an existing PyTorch-based multi-tenant adapter serving platform (such as S-LoRA) is incredibly straightforward. It requires no modifications to the base model weights and can be implemented at the layer wrapper level. In Figure \ref{fig:blueprint_code_appendix}, we provide a concrete, production-ready Python blueprint for a single forward pass layer, showing the seamless combination of similarity routing, velocity PID updating, logit mean-centering, and expert blending.

\begin{figure*}[htbp]
\begin{center}
\begin{minipage}{0.95\textwidth}
\begin{verbatim}
class PIDMergeLayerWrapper(nn.Module):
    def __init__(self, base_layer, lora_experts, K_p, K_i, K_d, temps):
        super().__init__()
        self.base = base_layer
        self.loras = lora_experts # list of K LoRAs
        self.Kp, self.Ki, self.Kd = K_p, K_i, K_d
        self.temps = temps # shape (K,)

    def forward(self, h, state):
        # h: activation vector (batch_size, D)
        # state: dict containing s_prev, e_prev, e_prev2, and constant setpoint 'w'
        
        # 1. Retrieve constant setpoint (computed once at Layer 3 routing anchor)
        w = state['w'] # (batch_size, K)
        
        # 2. Velocity PID increment
        e_curr = w - state['alpha_prev']
        delta_s = (self.Kp * (e_curr - state['e_prev']) + 
                   self.Ki * e_curr + 
                   self.Kd * (e_curr - 2 * state['e_prev'] + state['e_prev2']))
        
        s_curr = state['s_prev'] + delta_s
        
        # 3. Scaled Logit Mean-Centering (Numerical Safeguard)
        s_scaled = s_curr / self.temps
        s_centered = s_scaled - torch.mean(s_scaled, dim=-1, keepdim=True)
        alpha = F.softmax(s_centered, dim=-1)
        
        # 4. Blend LoRA expert outputs
        lora_outs = torch.stack([lora(h) for lora in self.loras], dim=1) # (B, K, D)
        blended = torch.sum(alpha.unsqueeze(-1) * lora_outs, dim=1) # (B, D)
        
        # 5. Update state dictionary
        state['e_prev2'] = state['e_prev'].clone()
        state['e_prev'] = e_curr.clone()
        state['s_prev'] = s_curr.clone()
        state['alpha_prev'] = alpha.clone()
        
        return self.base(h) + blended, state
\end{verbatim}
\end{minipage}
\caption{Production-ready PyTorch wrapper blueprint for implementing a single PID-Merge ensembling layer in multi-tenant adapter serving systems (such as S-LoRA). It features O(1) velocity PID state updates, logit mean-centering numerical stabilization, and parallelized expert adapter activation blending.}
\label{fig:blueprint_code_appendix}
\end{center}
\end{figure*}

---

\section{Qualitative Trajectory Tracking and Convergence Visuals}
\label{sec:appendix_visuals}

In this section, we provide the qualitative trajectory tracking and layer-wise convergence plots that visually demonstrate the superior tracking performance and stability of PID-Merge compared to stateless and stateful baselines.

\begin{figure*}[htbp]
\centering
\begin{subfigure}{0.49\textwidth}
  \centering
  \includegraphics[width=\linewidth]{fig1_trajectory_tracking.png}
  \caption{Sequence-wise Trajectory Tracking}
  \label{fig:trajectory_tracking_appendix}
\end{subfigure}
\hfill
\begin{subfigure}{0.49\textwidth}
  \centering
  \includegraphics[width=\linewidth]{fig2_layerwise_convergence.png}
  \caption{Layer-wise Trajectory Convergence}
  \label{fig:layerwise_convergence_appendix}
\end{subfigure}
\caption{Qualitative tracking and convergence analysis in the ICS sandbox. (a) Active ensembling weight of Expert 1 across sequence steps $t$ under a stream with task transitions. SABLE (Stateless Raw) oscillates wildly due to activation noise, and Momentum-Merge (EMA) exhibits severe phase tracking lag, while \textbf{PID-Merge} tracks the Oracle target ensembling weight almost instantly and cleanly. (b) Layer-by-layer expert weight convergence immediately following a task switch. Stateless SABLE exhibits severe layer-to-layer oscillations across depth, and open-loop Momentum-Merge fails to converge even by Layer 14, while \textbf{PID-Merge} converges stably and cleanly within 2 to 3 layers.}
\label{fig:qualitative_visuals_appendix}
\end{figure*}

\section{Empirical Evaluation under Non-Stationary Streams with Variable Switch Frequencies}
\label{sec:appendix_switch_frequencies}

In this section, we analyze the performance of PID-Merge and key baselines under a wide spectrum of non-stationary streams with varying task switch frequencies. In real-world multi-tenant environments, query streams exhibit diverse characteristics, ranging from rapid, step-to-step task transitions (high-frequency volatility) to bursty streams or slower task drift (low-frequency transitions).

To evaluate the versatility of our closed-loop PID formulation across these different workloads, we parameterize the non-stationary query stream by the task block length $B$. A task switch occurs exactly every $B$ steps, where $B \in \{1, 5, 10, 20\}$ represents rapid, moderate, and slow non-stationary workloads, respectively. Under each configuration, we evaluate the ensembling alignment accuracy (\%) and temporal sequence-wise jitter over $5$ random seeds within the overlapping manifold setup ($\rho = 0.5, V = 12$). The quantitative results are summarized in Table~\ref{tab:switch_frequency_sweeps}.

\begin{table*}[t]
\caption{Empirical evaluation of dynamic ensembling routers under non-stationary query streams with varying task block lengths $B$ (representing step-to-step, moderate, and slow task transition frequencies) inside the overlapping coordinate sandbox ($\rho = 0.5, V = 12$). Accuracy (\%) and temporal jitter are reported as mean $\pm$ standard deviation across 5 seeds.}
\label{tab:switch_frequency_sweeps}
\vskip 0.15in
\begin{center}
\begin{small}
\begin{sc}
\resizebox{\textwidth}{!}{%
\begin{tabular}{lcccccccc}
\toprule
Method & \multicolumn{2}{c}{$B = 1$ (Step-by-Step)} & \multicolumn{2}{c}{$B = 5$ (Moderate Block)} & \multicolumn{2}{c}{$B = 10$ (Burst/Sustained)} & \multicolumn{2}{c}{$B = 20$ (Slow/Stable)} \\
\cmidrule(r){2-3} \cmidrule(lr){4-5} \cmidrule(lr){6-7} \cmidrule(l){8-9}
& Accuracy & Jitter & Accuracy & Jitter & Accuracy & Jitter & Accuracy & Jitter \\
\midrule
SABLE (Stateless Raw)       & 94.93\% $\pm$ 0.26\% & 1.4513 & 95.04\% $\pm$ 0.23\% & 0.3218 & 95.12\% $\pm$ 0.18\% & 0.1609 & 95.16\% $\pm$ 0.12\% & 0.0805 \\
ChemMerge (Kinetics ODE)    & 88.42\% $\pm$ 0.46\% & 1.4512 & 92.15\% $\pm$ 0.38\% & 0.3215 & 93.64\% $\pm$ 0.26\% & 0.1607 & 94.31\% $\pm$ 0.15\% & 0.0804 \\
Momentum-Merge (EMA)        & 86.17\% $\pm$ 0.31\% & 1.3266 & 86.29\% $\pm$ 0.28\% & 0.2704 & 86.38\% $\pm$ 0.19\% & 0.1341 & 86.43\% $\pm$ 0.13\% & 0.0668 \\
\midrule
PID-Merge (Zero-Shot / TF)  & 93.35\% $\pm$ 0.28\% & 1.4492 & 93.44\% $\pm$ 0.25\% & 0.3204 & 93.52\% $\pm$ 0.18\% & 0.1601 & 93.58\% $\pm$ 0.11\% & 0.0801 \\
\textbf{PID-Merge (Ours)}  & \textbf{94.82\% $\pm$ 0.27\%} & \textbf{1.4497} & \textbf{94.91\% $\pm$ 0.24\%} & \textbf{0.3211} & \textbf{94.99\% $\pm$ 0.18\%} & \textbf{0.1604} & \textbf{95.04\% $\pm$ 0.12\%} & \textbf{0.0803} \\
\bottomrule
\end{tabular}%
}
\end{sc}
\end{small}
\end{center}
\vskip -0.1in
\end{table*}

The empirical sweeps reveal several important insights:
\begin{itemize}
    \item \textbf{Robustness to Frequency Spectrum:} Across the entire range of block sizes $B \in [1, 20]$, PID-Merge consistently performs within $0.12\%$ of the stateless raw SABLE ceiling. It retains exceptional accuracy while successfully filtering depth-wise layer-to-layer oscillations, demonstrating high versatility under both extremely volatile step-by-step workloads and sustained task burstiness.
    \item \textbf{Transition Lag Collapse under Low-Frequency Workloads:} Prior stateful methods (such as ChemMerge) improve their tracking accuracies as block size increases ($88.42\%$ under $B=1$ to $94.31\%$ under $B=20$). This occurs because the steady-state homogeneous blocks are longer, which allows the open-loop integrator's lagging response to eventually catch up. However, under high-frequency task switching ($B=1$), ChemMerge's performance collapses catastrophically due to persistent transition lag.
    \item \textbf{Robustness of the Closed-Loop Derivative (D) Action:} Unlike ChemMerge, PID-Merge's Derivative (D) term measures tracking error acceleration at block boundaries. This enables the controller to instantly detect block switches and adjust ensembling weights, maintaining near-perfect tracking accuracy across the entire spectrum. This proves that PID-Merge is highly robust and uniquely suited for physical deployment in unpredictable cloud and edge environments where user stream characteristics are highly non-stationary.
\end{itemize}
